% T5 fig:flux (opus3) — structure reinvention: the two-line crossing is generalized to a FAMILY
% over affinity A. Reverse flux r^- xi (fixed, r^-=1) is crossed by three forward lines
% r^+(1-xi) for r^+ in {0.5,1,2} = exp(A/RT) in {A<0, A=0, A>0}; the stationary point
% xi_eq = r^+/(r^+ + r^-) slides 1/3 -> 1/2 -> 2/3 along the diagonal. All values from T5_calc.py.
\centering
\begin{tikzpicture}[x=4.0cm,y=1.15cm]
\draw[-{Latex[length=1.4mm]}] (0,0) -- (0,2.35) node[above,font=\scriptsize] {flux [arb.]};
\draw[-{Latex[length=1.4mm]}] (-0.01,0) -- (1.12,0) node[below,font=\scriptsize] {$\xi$};
\foreach \xx in {0.5,1} {\draw (\xx,0.03) -- (\xx,-0.03) node[below,font=\scriptsize] {\xx};}
% reverse flux (fixed): r- xi , the diagonal that all crossings sit on
\draw[semithick] (0,0) -- (1,1);
\node[right,font=\scriptsize] at (0.80,0.66) {reverse $r^-\xi$};
% forward lines r+(1-xi)
\draw[very thick] (0,2) -- (1,0);                    % A>0, r+=2
\draw[gray,semithick] (0,1) -- (1,0);                % A=0, r+=1
\draw[densely dashed,semithick] (0,0.5) -- (1,0);    % A<0, r+=0.5
\node[anchor=south west,font=\scriptsize] at (0.04,1.93) {forward $r^+(1{-}\xi)$};
% affinity legend (upper-right clear zone)
\node[anchor=west,font=\tiny] at (0.66,1.92) {$\mathcal A{>}0$: $r^+{=}2$};
\node[anchor=west,font=\tiny,gray] at (0.66,1.70) {$\mathcal A{=}0$: $r^+{=}1$};
\node[anchor=west,font=\tiny] at (0.66,1.48) {$\mathcal A{<}0$: $r^+{=}0.5$};
% crossings (all on the diagonal) — xi_eq slides right as A grows
\fill (0.6667,0.6667) circle (1.0pt) node[above left,font=\scriptsize,inner sep=1.5pt] {$\tfrac23$};
\fill (0.5,0.5) circle (1.0pt) node[above left,font=\tiny,inner sep=1.5pt] {$\tfrac12$};
\fill (0.3333,0.3333) circle (1.0pt) node[below right,font=\scriptsize,inner sep=1.5pt] {$\tfrac13$};
\draw[densely dotted] (0.6667,0)--(0.6667,0.6667);
\draw[densely dotted] (0.3333,0)--(0.3333,0.3333);
\node[below,font=\scriptsize] at (0.5,-0.42) {$\xi_\eq=\dfrac{r^+}{r^++r^-}=\dfrac{1}{1+e^{-\mathcal A/RT}}$ (식~\eqref{eq:logisticsolve})};
\end{tikzpicture}
\caption{정지점 유도 --- 친화도 $\mathcal A$ 에 대한 \emph{족}으로 재구성했다. 고정된 역방향 플럭스
$r^-\xi$(대각선)를 세 정방향 직선 $r^+(1-\xi)$ 이 가로지르며, 각 교점이 정지점 $\xi_\eq=r^+/(r^++r^-)$
(식~\eqref{eq:logisticsolve})이다. detailed balance $r^+/r^-=e^{\mathcal A/RT}$(식~\eqref{eq:db})가
$r^+\in\{0.5,1,2\}$ 를 정해 $\xi_\eq$ 가 $\tfrac13\!\to\!\tfrac12\!\to\!\tfrac23$ 으로 오른쪽으로 이동한다
(좌표는 식 그대로의 수치 평가). 세 교점이 모두 역방향 대각선 위에 놓이는 것이 $\xi_\eq$ 가 logistic
$1/(1+e^{-\mathcal A/RT})$ 과 같음을 시각적으로 보인다.}
\label{fig:flux}
